\documentclass{article}
\usepackage[utf8]{inputenc}
\usepackage{verbatim}
\usepackage{siunitx}
\usepackage{listings}
\usepackage{xcolor}

\usepackage[T1]{fontenc}
\usepackage{fontspec}


\definecolor{codegreen}{rgb}{0,0.6,0}
\definecolor{codegray}{rgb}{0.5,0.5,0.5}
\definecolor{codepurple}{rgb}{0.58,0,0.82}
\definecolor{backcolour}{rgb}{0.95,0.95,0.92}

\definecolor{linkblue}{rgb}{0.1,0.3,0.7}

\lstdefinestyle{mystyle}{
    backgroundcolor=\color{backcolour},   
    commentstyle=\color{codegreen},
    keywordstyle=\color{magenta},
    numberstyle=\tiny\color{codegray},
    stringstyle=\color{codepurple},
    basicstyle=\ttfamily\footnotesize,
    breakatwhitespace=false,         
    breaklines=true,                 
    captionpos=b,                    
    keepspaces=true,                 
    numbers=left,                    
    numbersep=5pt,                  
    showspaces=false,                
    showstringspaces=false,
    showtabs=false,                  
    tabsize=2
}

\lstset{style=mystyle}



\usepackage{setspace}
%\usepackage{pgfplots}
\usepackage{tikz}
\usepackage{mathtools}
\usetikzlibrary{arrows.meta}

\onehalfspacing % better to read

\usepackage{multirow} %multiple row for tabular env.
\usepackage{tabularx} %table
\usepackage{makecell}
\usepackage{booktabs}
\usepackage{amsmath,amssymb,amstext} %formulas ect.
\newcommand{\angstrom}{\textup{\AA}} %the unit
\usepackage{indentfirst} %indest first paragraph

\usepackage{csquotes}

\usepackage[hidelinks]{hyperref} %Links to labeled formulas
\hypersetup{
    colorlinks=true,
    linkcolor=linkblue,
    filecolor=magenta,      
    urlcolor=linkblue,
    citecolor=linkblue,
    pdfpagemode=FullScreen,
    }

\usepackage[edges]{forest}
\usepackage{color, colortbl}
\usepackage{geometry} % more space to write on page
\setlength{\topmargin}{-2.5cm}
\setlength{\oddsidemargin}{0cm}
\setlength{\evensidemargin}{0cm}
\setlength{\textwidth}{17cm}
\setlength{\textheight}{23cm}

\setstretch{1.15}
\setlength{\parskip}{6em plus 0.2em}  % Space between paragraphs

\usepackage{graphicx} % Subfigures
\usepackage{subcaption}

%\usepackage{natbib}
\usepackage{graphicx}
\usepackage{tikz}
\usetikzlibrary{arrows}
\usetikzlibrary{calc}
\usetikzlibrary{angles,quotes}
\usepackage{centernot}

%%%%%%%%%%new stuff i added for bachelors
\newcommand{\draft}{\textcolor{red}{Draft:}}
\newcommand{\temp}{\textcolor{red}{Temp:}}
\newcommand{\source}{\textcolor{red}{[source?]}}
\newcommand{\hm}{\textcolor{red}{[hmm?]}}
\newcommand{\form}{\textcolor{red}{[formulation]}}

%\newcommand{\note}[1]{\textcolor{red}{#1}}
\usepackage{gensymb}
\usepackage{tkz-euclide}

\usepackage{ifthen}
\newcommand{\note}[1]{%
    \ifthenelse{\boolean{show_notes}}{%
        \textcolor{red}{#1}%
    }{}%
}

%add numbering thing on equation manually
\newcommand\numberthis{\addtocounter{equation}{1}\tag{\theequation}}

%\setcitestyle{square}

%set darkmode
%\pagecolor[rgb]{0,0,0}
%\color[rgb]{0.7,0.7,0.7}


\graphicspath{ {./} }


\begin{document}


\newboolean{show_notes}

\setboolean{show_notes}{false}




\begin{titlepage}
\begin{center}
\vspace*{4cm}
\huge
{\fontfamily{lmtt} Comparing Algorithms Algorithming the Maximum Algorithms
Comparing Problem}
\vspace*{1cm}
\large

by Author

student number: 123

E-Mail: Mail 
\vspace*{1cm}


Bachelor Thesis

Bachelor of Science

Major in Computer Science
\vspace*{1cm}

First Supervisor: 

Second Supervisor: 

University
\vspace*{1cm}

Date of Submission: Submission date
\end{center}
\end{titlepage}


\newpage


\vspace*{2cm}
\begin{center}
\begin{minipage}{0.7\textwidth}
{\large \hspace{2pt}\textbf{Abstract}\vspace{1pt}}\\

The approximation of the Travelling Salesperson Problem is a widely studied field
where even simple approaches seem to promise good results. In this thesis, we
examine two different approaches to approximating the Maximum TSP. The first, a
Graph Patch Heuristic, attempts to find a solution by patching a maximum-weight
cycle cover into a singular cycle. The second, called the Tunnelling Algorithm,
builds on the idea of approximating the Euclidean norm with a metric defined by
a polytope. Despite their drastically different methodologies, we find that both
approaches reduce to a similar sub-problem and yield similar approximation
ratios. While the Graph Patch Heuristic seems to outperform the Tunnelling
Algorithm with larger input sizes and dimensions, we find that the Tunnelling
Algorithm actually solves the Max TSP in polynomial time for common norms such
as $L_1$ and $L_\infty$.

This thesis expands on some of the nuances of both algorithms, presenting them
in a self-contained manner.

\end{minipage}
\end{center}

\newpage
{
\hypersetup{linkcolor=black}
\tableofcontents
}
\newpage

\include{src/list}



\section{Introduction}
\label{sec:intro}

% NP

Since the relationship between the set of Problems in $P$, solvable in
polynomial time, and Problems in $NP$, the set of problems where the solutions
can be verified in polynomial time, is unclear, there exists great interest in
finding approximation algorithms solving NP and NP-hard problems. 

%explaining approximations

An approximation algorithm is supposed to guarantee a better runtime, with the
trade-off that the solution may not be optimal. Therefore, there are at least
two factors to compare approximation algorithms by, runtime and how close the
solution is to the optimal solution. There are multiple ways to evaluate the
accuracy of an algorithm, in some cases it can be an additive factor, but most
commonly, the accuracy is measured with an approximation ratio $\rho(n)$, where
the value of a solution $Approx$ provided by an approximation algorithm is
within a factor of $\rho(n)$ to the actual optimal value $OPT$ for any input
size $n$. Conventionally $\rho(n)\geq 1$ by defining it as $\frac{OPT}{Approx}$
for maximization problems and $\frac{Approx}{OPT}$ for minimization problems, or
generalizing to $max(\frac{OPT}{Approx}, \frac{Approx}{OPT})\leq
\rho(n)${algorithms}. Some authors, such as
Shenmaier{shenmaier1} have used a relative error $\delta$, where $\delta =
\frac{ OPT-Approx}{OPT}$ to describe the accuracy of the approximation
algorithm. The APX class is a category of algorithms that approximate an
optimisation problem in NP in polynomial runtime with an approximation ratio
that is constant. 

\section{Greedy Patch Heuristic}
\label{sec:gph}

We will start with the simpler one of the two algorithms that finds a tour by
first finding a maximum-weight cycle cover and then combining it into one cycle
creating a tour. The results of this section are mainly focused on the paper by
Shenmaier(2022){shenmaier1}, where he describes an algorithm for which he
estimates a constant approximation ratio, but the more interesting result is
that he proves a relative error that improves with input size. The algorithm
itself is a modified version of the greedy Karp-Steel algorithm where
Shenmaier{shenmaier1} references papers that analyse such greedy
Karp-Steel algorithm in the context of Minimum TSP. For more exact steps of the
algorithm we will mainly use details from a 1999 paper by Clover et
al.{karpsteel}, where we will of course modify the proposed algorithm to
work for the Maximum TSP. Additionally the Hungarian Algorithm from \[...\] will find
its use in this algorithm. Tracking down the references of various papers
discussing such a patching algorithm we can find that the idea dates back to at
leasts a 1976 paper by Karp{karp}, as well as another paper we consider
from 1985 by Karp and Steel{oldkarpsteel}.

\end{document}

